
Three recent LLM formal reasoning systems share a common design feature: they supply the language model with structured feedback from a formal verifier. BFS-Prover achieves 72.95\% on the miniF2F theorem proving benchmark using Lean4 compiler errors as DPO negatives \citep{Xinetal2025}. PropertyGPT improves smart contract property compilation success from 63\% to 87\% using static analysis feedback at inference time \citep{Liuetal2024}. Proof of Thought reduces compilation errors from 14.6\% to 0\% using Z3 SMT counterexamples via a JSON DSL bridge \citep{Gangulyetal2024}. None of these systems identifies the mechanism by which formal feedback produces these improvements.

Two mechanistically distinct explanations are consistent with the observed results. Under the \textbf{oracle hypothesis}, structured formal feedback carries state-specific semantic information that directs the policy toward formally valid continuations. Under the \textbf{regularizer hypothesis}, formal feedback acts as a source of contrastive examples for DPO training, sharpening the policy through a general regularization effect that does not depend on the semantic content of the negative example.

These hypotheses predict different design conclusions. If the oracle hypothesis holds, feedback quality (semantic specificity) is the critical design variable. If the regularizer hypothesis holds, feedback diversity (number and variety of negatives) is more important. These conclusions point in opposite directions for future system design.

No existing work has been designed to distinguish these hypotheses. Each system uses a single feedback type; cross-condition controlled ablations have not been performed.

This paper introduces a 2×2 factorial design within the BFS-Prover framework intended to test the oracle/regularizer distinction. The design holds all variables constant except the semantic content of the DPO negative example, using three conditions: step-local Lean4 compiler errors (grounded), same-state failed-branch tactics (ungrounded), and permuted compiler error messages (control). The \textbf{locality score} metric is introduced as a mechanistic probe measuring what fraction of the post-DPO probability mass shift concentrates on tactic categories consistent with the specific error at each failed proof state.

Attempting to execute this experiment produced an infrastructure failure rather than empirical results. The LeanDojo proof state extraction layer silently fell back to synthetic data generation when the Lean4 toolchain was absent, producing uniformly zero locality scores that passed format validation while carrying no scientific content. This failure mode and the proposed remediation protocol are documented as a methodological contribution.

\textbf{Contributions:}

\begin{enumerate}
\item The \textbf{oracle/regularizer framing}: a testable mechanistic distinction that unifies BFS-Prover, PropertyGPT, and Proof of Thought under a single causal hypothesis.
\end{enumerate}

\begin{enumerate}
\item The \textbf{locality score metric}: a novel mechanistic probe — the fraction of post-DPO probability mass shift on premise-consistent tactic categories — that directly operationalizes oracle function versus diffuse regularization.
\end{enumerate}

\begin{enumerate}
\item An \textbf{implemented and unit-tested experimental infrastructure}: a 3-condition DPO pipeline with 100\% state alignment verification, pre-specified tactic taxonomy (type_error, undefined_name, tactic_failure), and correct DPO loss implementation (β=10, lr 5e-6→5e-7, average_log_prob=False), ready for re-execution with real LeanDojo data.
\end{enumerate}

\begin{enumerate}
\item A \textbf{methodological warning}: documentation of the silent synthetic fallback failure mode in LLM+formal-verifier pipelines, with a proposed pre-run environment validation protocol.
\end{enumerate}
